\documentclass[12pt,a4paper]{article}
\usepackage[utf8]{inputenc}
\usepackage[margin=1in]{geometry}
\usepackage{hyperref}
\usepackage{graphicx}
\usepackage{listings}
\usepackage{xcolor}
\usepackage{enumitem}
\usepackage{tabularx}
\usepackage{booktabs}

% Code listing settings
\lstset{
    basicstyle=\ttfamily\small,
    breaklines=true,
    frame=single,
    backgroundcolor=\color{gray!10}
}

\title{\textbf{NeuroBloom\\Cognitive Training Platform}\\System Documentation}
\author{}
\date{January 2026}

\begin{document}

\maketitle
\tableofcontents
\newpage

\section{Overview}

\textbf{NeuroBloom} is a comprehensive cognitive training platform designed to assess and improve cognitive abilities across six key domains. The system uses adaptive algorithms to personalize training difficulty and provides detailed performance analytics.

\subsection{Core Features}

\begin{itemize}
    \item User authentication (Register/Login)
    \item Baseline cognitive assessment (6 domains)
    \item Personalized training plan generation
    \item Adaptive difficulty adjustment
    \item Performance tracking and visualization
    \item Badge achievement system
    \item Daily streak tracking with freeze mechanism
    \item Real-time progress monitoring
    \item Baseline vs Current performance comparison
\end{itemize}

\section{System Architecture}

\subsection{Backend}
\begin{itemize}
    \item \textbf{Framework}: FastAPI (Python)
    \item \textbf{Database}: PostgreSQL
    \item \textbf{ORM}: SQLModel
    \item \textbf{Authentication}: Token-based
\end{itemize}

\subsection{Frontend}
\begin{itemize}
    \item \textbf{Framework}: SvelteKit (Svelte 5)
    \item \textbf{Language}: TypeScript
    \item \textbf{Styling}: Custom CSS with gradients
    \item \textbf{API Communication}: Axios
\end{itemize}

\subsection{Key Models}
\begin{enumerate}
    \item \textbf{User} -- Authentication and profile data
    \item \textbf{BaselineAssessment} -- Initial cognitive profile
    \item \textbf{TrainingPlan} -- Personalized training configuration
    \item \textbf{TrainingSession} -- Individual task completion records
    \item \textbf{CognitiveTask} -- Task definitions and metadata
    \item \textbf{Badge} -- Achievement tracking
\end{enumerate}

\section{User Journey}

\subsection{Registration and Login}

\begin{center}
Register $\rightarrow$ Login $\rightarrow$ Dashboard
\end{center}

\begin{itemize}
    \item \textbf{Registration}: Email + Password
    \item \textbf{Login}: Token-based authentication
    \item \textbf{Session Management}: Persistent login state
\end{itemize}

\subsection{Baseline Assessment}

\begin{center}
Dashboard $\rightarrow$ Complete 6 Tasks $\rightarrow$ Calculate Baseline $\rightarrow$ Training Plan
\end{center}

\subsubsection{Six Baseline Tasks}
\begin{enumerate}
    \item \textbf{Working Memory} -- N-back task
    \item \textbf{Processing Speed} -- Simple Reaction Time
    \item \textbf{Attention} -- Flanker task
    \item \textbf{Cognitive Flexibility} -- Task Switching
    \item \textbf{Planning} -- Tower of London
    \item \textbf{Visual Scanning} -- Visual Search
\end{enumerate}

\subsubsection{Assessment Flow}
\begin{enumerate}
    \item User completes all 6 tasks (one per domain)
    \item System calculates domain scores (0-100)
    \item Overall cognitive score computed (average of 6 domains)
    \item Baseline results page displays comprehensive analysis
    \item Training plan generated based on baseline scores
\end{enumerate}

\subsection{Training Phase}

\begin{center}
Training Plan $\rightarrow$ Complete Sessions $\rightarrow$ Track Progress $\rightarrow$ View Analytics
\end{center}

\subsubsection{Session Structure}
\begin{itemize}
    \item Each session = \textbf{4 tasks}
    \item Tasks selected from different domains
    \item Task order randomized for variety
    \item Progress tracked: ``1/4'', ``2/4'', ``3/4'', ``4/4''
\end{itemize}

\subsection{Progress Monitoring}

\begin{center}
Dashboard $\rightarrow$ Training Page $\rightarrow$ Progress Page $\rightarrow$ Baseline Results
\end{center}

\section{Cognitive Domains}

\subsection{Working Memory}

\textbf{Definition:} Ability to hold and manipulate information in mind temporarily.

\textbf{Tasks:}
\begin{itemize}
    \item N-Back
    \item Digit Span
    \item Spatial Span
    \item Letter-Number Sequencing
    \item Operation Span
\end{itemize}

\textbf{Real-world applications:} Remembering phone numbers, following multi-step instructions, mental arithmetic

\subsection{Attention}

\textbf{Definition:} Ability to focus on relevant information while ignoring distractions.

\textbf{Tasks:}
\begin{itemize}
    \item Flanker Task
    \item Go/No-Go
    \item Stroop Test
\end{itemize}

\textbf{Real-world applications:} Driving in traffic, studying in noisy environments, selective listening

\subsection{Cognitive Flexibility}

\textbf{Definition:} Ability to switch between different tasks or mental sets.

\textbf{Tasks:}
\begin{itemize}
    \item Task Switching
    \item DCCS (Dimensional Change Card Sort)
    \item Plus-Minus Task
\end{itemize}

\textbf{Real-world applications:} Multitasking, adapting to new situations, problem-solving from different angles

\subsection{Planning}

\textbf{Definition:} Ability to organize steps to achieve a goal.

\textbf{Tasks:}
\begin{itemize}
    \item Tower of London
    \item Twenty Questions
    \item Category Fluency
\end{itemize}

\textbf{Real-world applications:} Project management, trip planning, organizing daily schedules

\subsection{Processing Speed}

\textbf{Definition:} How quickly you can process and respond to information.

\textbf{Tasks:}
\begin{itemize}
    \item Simple Reaction Time
    \item SDMT (Symbol Digit Modalities Test)
    \item Trail Making Test (Part A)
    \item Inspection Time
    \item Pattern Comparison
\end{itemize}

\textbf{Real-world applications:} Quick decision-making, sports performance, typing speed

\subsection{Visual Scanning}

\textbf{Definition:} Ability to quickly locate specific information in a visual field.

\textbf{Tasks:}
\begin{itemize}
    \item Visual Search
    \item Cancellation Test
    \item Multiple Object Tracking
\end{itemize}

\textbf{Real-world applications:} Reading maps, finding items in a cluttered space, visual proofreading

\section{Training System}

\subsection{Focus Area Classification}

Based on baseline scores, domains are categorized into:

\begin{enumerate}
    \item \textbf{Primary Focus (2 domains)}
    \begin{itemize}
        \item Weakest areas
        \item Highest training priority
        \item More frequent task selection
    \end{itemize}
    
    \item \textbf{Secondary Focus (2 domains)}
    \begin{itemize}
        \item Middle-tier performance
        \item Moderate training frequency
        \item Balanced improvement
    \end{itemize}
    
    \item \textbf{Maintenance (2 domains)}
    \begin{itemize}
        \item Strongest areas
        \item Periodic practice
        \item Prevent skill degradation
    \end{itemize}
\end{enumerate}

\subsection{Task Selection Algorithm}

Each training session selects 4 tasks:
\begin{enumerate}
    \item Random domain selection (varied each session)
    \item Task rotation within domain (cycling through available tasks)
    \item Difficulty level based on current performance
    \item No duplicate domains in single session
\end{enumerate}

\section{Adaptive Difficulty Algorithm}

\subsection{How It Works}

After each session (4 tasks complete):
\begin{enumerate}
    \item Evaluate performance per domain
    \item Adjust difficulty based on accuracy
    \item Update difficulty levels in training plan
\end{enumerate}

\subsection{Difficulty Levels}

\begin{itemize}
    \item \textbf{Range}: 1 (easiest) to 10 (hardest)
    \item \textbf{Starting point}: Based on baseline score
    \begin{itemize}
        \item Score < 50: Difficulty 3
        \item Score 50-70: Difficulty 5
        \item Score > 70: Difficulty 7
    \end{itemize}
\end{itemize}

\subsection{Adjustment Rules}

\begin{lstlisting}[language=Python]
if accuracy >= 85%:
    difficulty = min(current + 1, 10)  # Increase (max 10)
elif accuracy < 65%:
    difficulty = max(current - 1, 1)   # Decrease (min 1)
else:
    difficulty = current               # Maintain
\end{lstlisting}

\textbf{Logic:}
\begin{itemize}
    \item \textbf{High accuracy ($\geq$85\%)} $\rightarrow$ Task too easy $\rightarrow$ Increase difficulty
    \item \textbf{Low accuracy (<65\%)} $\rightarrow$ Task too hard $\rightarrow$ Decrease difficulty
    \item \textbf{Medium accuracy (65-85\%)} $\rightarrow$ Optimal challenge $\rightarrow$ Maintain
\end{itemize}

\subsection{Difficulty Updates}

\begin{itemize}
    \item Updates happen \textbf{only after full session} (4 tasks complete)
    \item Each domain adjusted independently
    \item Changes persist across sessions
    \item Real-time display on Training page
\end{itemize}

\section{Session Structure}

\subsection{Session Lifecycle}

\subsubsection{Session Start}
\begin{itemize}
    \item 4 tasks assigned (random domains)
    \item Difficulty set per domain
    \item Progress: 0/4
\end{itemize}

\subsubsection{Task Completion}
\begin{itemize}
    \item User completes task
    \item Score calculated
    \item Progress updated (1/4, 2/4, 3/4)
    \item Redirects to Training page
\end{itemize}

\subsubsection{Session Complete (4/4 tasks done)}
\begin{itemize}
    \item Difficulty adjustment calculated
    \item Streak updated
    \item Badges checked
    \item Session counter incremented
    \item New session initialized
\end{itemize}

\subsection{Key Metrics Per Task}

\begin{itemize}
    \item \textbf{Score}: Overall performance (0-100)
    \item \textbf{Accuracy}: Percentage correct responses
    \item \textbf{Reaction Time}: Average response speed (ms)
    \item \textbf{Consistency}: Standard deviation of reaction times
    \item \textbf{Errors}: Number of mistakes
\end{itemize}

\section{Scoring and Metrics}

\subsection{Score Calculation}

\subsubsection{Baseline Scores}
\begin{itemize}
    \item Calculated from 6 baseline tasks
    \item Domain scores: 0-100 scale
    \item Overall score: Average of 6 domains
\end{itemize}

\subsubsection{Training Scores}
\begin{itemize}
    \item Each task scored 0-100
    \item Domain average: Mean of all tasks in that domain
    \item Overall average: Mean across all domains
\end{itemize}

\subsection{Score Labels}

\begin{table}[h]
\centering
\begin{tabular}{ll}
\toprule
\textbf{Score Range} & \textbf{Label} \\
\midrule
$\geq$ 80 & Excellent (Green) \\
70-79 & Good (Light Green) \\
60-69 & Average (Orange) \\
50-59 & Below Average (Light Red) \\
< 50 & Needs Improvement (Red) \\
\bottomrule
\end{tabular}
\end{table}

\subsection{Dashboard Metrics}

\begin{enumerate}
    \item \textbf{Total Sessions}
    \begin{itemize}
        \item Count of completed 4-task sessions
        \item NOT individual task count
    \end{itemize}
    
    \item \textbf{Total Tasks}
    \begin{itemize}
        \item Count of individual tasks completed
        \item 4 tasks = 1 session
    \end{itemize}
    
    \item \textbf{Active Domains}
    \begin{itemize}
        \item Number of domains with training history
        \item Maximum: 6 domains
    \end{itemize}
    
    \item \textbf{Last Training}
    \begin{itemize}
        \item Date of most recent session
    \end{itemize}
\end{enumerate}

\section{Visualizations and Graphs}

\subsection{Radar Chart (Cognitive Profile)}

\textbf{What it shows:}
\begin{itemize}
    \item Hexagonal chart with 6 axes (one per domain)
    \item Each axis represents one cognitive domain
    \item Distance from center = performance score
\end{itemize}

\textbf{Two overlays:}
\begin{itemize}
    \item \textbf{Baseline (Blue dashed line)}: Initial assessment scores
    \item \textbf{Current (Green solid line)}: Current training performance
\end{itemize}

\textbf{How to read:}
\begin{itemize}
    \item Larger area = Better overall performance
    \item Symmetrical shape = Balanced cognitive profile
    \item Asymmetrical = Strength/weakness pattern
    \item Gap between lines = Improvement/decline
\end{itemize}

\textbf{Purpose:}
\begin{itemize}
    \item Quick visual comparison
    \item Identify strongest/weakest domains
    \item Track progress over time
\end{itemize}

\subsection{Domain Score Cards}

\textbf{Display format:}
\begin{verbatim}
┌─────────────────────────┐
│  Working Memory         │
│                         │
│  [72.5]     [78.3]     │
│  Baseline   Current     │
└─────────────────────────┘
\end{verbatim}

\textbf{Features:}
\begin{itemize}
    \item Side-by-side comparison
    \item Color-coded by performance level
    \item Shows both baseline and current scores
\end{itemize}

\subsection{Progress Bars}

\textbf{Used for:}
\begin{itemize}
    \item Baseline task completion (6/6 tasks)
    \item Session progress (1/4, 2/4, 3/4, 4/4)
    \item Domain difficulty levels (visual bars)
\end{itemize}

\textbf{Visual feedback:}
\begin{itemize}
    \item Filled portion = Progress percentage
    \item Color = Performance level
    \item Animation on updates
\end{itemize}

\subsection{Difficulty Visualization}

\textbf{Training Page - Current Difficulty:}
\begin{verbatim}
Working Memory    [████████░░] Level 8/10
Attention         [██████░░░░] Level 6/10
Flexibility       [██░░░░░░░░] Level 2/10
Planning          [███████░░░] Level 7/10
Processing Speed  [█████████░] Level 9/10
Visual Scanning   [████░░░░░░] Level 4/10
\end{verbatim}

\textbf{Features:}
\begin{itemize}
    \item Real-time difficulty display
    \item Visual bar representation
    \item Numeric level (1-10)
\end{itemize}

\subsection{Overall Score Comparison}

\textbf{Two-card layout:}

\textbf{Baseline Score Card:}
\begin{itemize}
    \item Initial assessment date
    \item Overall score (average)
    \item Performance label
\end{itemize}

\textbf{Current Score Card:}
\begin{itemize}
    \item Latest performance
    \item Current overall score
    \item Improvement indicator (e.g., $\uparrow$ 5.2 points)
\end{itemize}

\section{Gamification Features}

\subsection{Badge System}

\subsubsection{Badge Categories}

\textbf{A. Streak Badges}
\begin{itemize}
    \item Getting Started -- 3 day streak
    \item On Fire -- 7 day streak
    \item Unstoppable -- 14 day streak
    \item Legendary -- 30 day streak
\end{itemize}

\textbf{B. Session Completion Badges}
\begin{itemize}
    \item First Steps -- 1 session
    \item Getting Going -- 5 sessions
    \item Committed -- 10 sessions
    \item Dedicated -- 25 sessions
    \item Expert -- 50 sessions
    \item Master -- 100 sessions
\end{itemize}

\textbf{C. Domain Mastery Badges}
\begin{itemize}
    \item Working Memory Master
    \item Attention Master
    \item Flexibility Master
    \item Planning Master
    \item Speed Master
    \item Visual Scanning Master
\end{itemize}

\textbf{Earning Criteria:}
\begin{itemize}
    \item Average score $\geq$ 80 in domain
    \item Minimum 10 tasks completed in domain
\end{itemize}

\textbf{Display:}
\begin{itemize}
    \item Notification on earning
    \item Badge gallery view
    \item Progress tracking
\end{itemize}

\subsection{Streak Tracking}

\textbf{Streak System:}
\begin{itemize}
    \item Tracks consecutive days trained
    \item Encourages daily practice
    \item Includes freeze mechanism
\end{itemize}

\textbf{How it works:}

\begin{verbatim}
Day 1: First session → Streak = 1
Day 2: Train → Streak = 2
Day 3: Train → Streak = 3
Day 4: SKIP → Freeze used (if available) → Streak = 3
Day 6: Train → Streak = 4
Day 8: SKIP (no freeze) → Streak = 1 (reset)
\end{verbatim}

\textbf{Streak Freeze:}
\begin{itemize}
    \item \textbf{One-time protection} per streak
    \item Protects against 1 missed day
    \item Replenishes when streak breaks
    \item Visual indicator on dashboard
\end{itemize}

\textbf{Streak Metrics:}
\begin{itemize}
    \item \textbf{Current Streak}: Consecutive days
    \item \textbf{Longest Streak}: Personal best
    \item \textbf{Total Training Days}: Lifetime count
\end{itemize}

\section{Performance Tracking}

\subsection{Performance Comparison API}

\textbf{Endpoint:} \texttt{/api/training/training-session/performance-comparison/\{user\_id\}}

\textbf{Returns:}
\begin{lstlisting}[language=json]
{
  "comparison": {
    "working_memory": {
      "baseline": 65.5,
      "current": 72.3,
      "change": 6.8,
      "tasks_completed": 12
    },
    "attention": { ... },
    ...
  }
}
\end{lstlisting}

\textbf{Features:}
\begin{itemize}
    \item Baseline vs current comparison
    \item Absolute change calculation
    \item Task count per domain
    \item Only includes domains with training data
\end{itemize}

\subsection{Metrics Dashboard}

\textbf{Endpoint:} \texttt{/api/training/training-session/metrics/\{user\_id\}}

\textbf{Returns:}
\begin{lstlisting}[language=json]
{
  "total_sessions": 5,
  "total_tasks": 20,
  "metrics_by_domain": {
    "working_memory": {
      "total_sessions": 5,
      "average_score": 72.3,
      "average_accuracy": 78.5,
      "current_difficulty": 6,
      "improvement": 8.2,
      "trend": "improving"
    }
  },
  "last_training_date": "2026-01-09T10:30:00"
}
\end{lstlisting}

\textbf{Trend Classification:}
\begin{itemize}
    \item \textbf{Improving}: Change > +5 points
    \item \textbf{Stable}: Change between -5 and +5
    \item \textbf{Declining}: Change < -5 points
\end{itemize}

\subsection{Session History}

\textbf{Tracking:}
\begin{itemize}
    \item Every task completion saved
    \item Timestamp recorded
    \item Performance metrics stored
    \item Difficulty levels logged
\end{itemize}

\textbf{Data retention:}
\begin{itemize}
    \item Unlimited history
    \item Accessible via API
    \item Used for trend analysis
\end{itemize}

\section{Technical Implementation}

\subsection{Cache Busting}

\textbf{Problem:} Browser caching stale training data

\textbf{Solution:} Timestamp query parameters
\begin{lstlisting}[language=JavaScript]
getPlan: async (userId) => {
  const response = await api.get(
    `/api/training/training-plan/${userId}?_t=${Date.now()}`
  );
  return response.data;
}
\end{lstlisting}

\textbf{Applied to:}
\begin{itemize}
    \item Training plan retrieval
    \item Metrics fetching
    \item Performance comparison
    \item Next tasks recommendation
\end{itemize}

\subsection{Session Completion Tracking}

\textbf{Implementation:}
\begin{lstlisting}[language=Python]
completed_tasks = plan.get_current_session_tasks_completed()
task_identifier = task_id if task_id else domain
if task_identifier not in completed_tasks:
    completed_tasks.append(task_identifier)
    plan.current_session_tasks_completed = json.dumps(
        completed_tasks
    )

session_complete = len(completed_tasks) >= 4

if session_complete:
    # Adjust difficulties
    # Update streak
    # Check badges
    # Reset session
    plan.total_sessions_completed += 1
    plan.current_session_tasks_completed = "[]"
\end{lstlisting}

\subsection{Difficulty Persistence}

\textbf{Storage:}
\begin{lstlisting}[language=Python]
# JSON format in database
current_difficulty = {
    "working_memory": 6,
    "attention": 5,
    "flexibility": 3,
    "planning": 7,
    "processing_speed": 8,
    "visual_scanning": 4
}
\end{lstlisting}

\textbf{Updates:}
\begin{itemize}
    \item Only after full session (4 tasks)
    \item Per-domain independent adjustment
    \item Persisted to database
    \item Retrieved for next session
\end{itemize}

\subsection{Baseline Recalculation}

\textbf{Feature:} Allow users to retake baseline assessment

\textbf{Implementation:}
\begin{enumerate}
    \item User clicks ``Recalculate Baseline''
    \item System fetches latest 6 baseline tasks
    \item Recalculates domain scores
    \item Updates \texttt{assessment\_date} timestamp
    \item Refreshes training plan if needed
\end{enumerate}

\textbf{Use cases:}
\begin{itemize}
    \item User improved and wants new baseline
    \item Initial assessment on bad day
    \item Progress milestone check
\end{itemize}

\section{Data Flow Diagrams}

\subsection{Task Completion Flow}

\begin{verbatim}
User completes task
    ↓
Submit results to backend
    ↓
Calculate score, accuracy, metrics
    ↓
Create TrainingSession record
    ↓
Add task_id to completed tasks array
    ↓
Check if session complete (4/4)
    ↓
YES → Adjust difficulties, update streak, check badges
    ↓
Redirect to Training page
    ↓
Page reloads fresh data (cache-busted)
    ↓
Display updated progress & difficulty
\end{verbatim}

\subsection{Streak Update Flow}

\begin{verbatim}
Session completes
    ↓
Get last_session_date from plan
    ↓
Calculate days_since_last
    ↓
If 0 days → Update timestamp, no streak change
If 1 day → Increment streak
If 2 days + freeze available → Use freeze, maintain streak
If 2+ days OR no freeze → Reset streak to 1
    ↓
Update longest_streak if current > longest
    ↓
Set last_session_date to now
\end{verbatim}

\section{Key Features Summary}

\subsection{Authentication System}
\begin{itemize}
    \item Secure registration and login
    \item Token-based sessions
    \item Persistent user state
\end{itemize}

\subsection{Baseline Assessment}
\begin{itemize}
    \item 6 domain evaluation
    \item Comprehensive scoring
    \item Visual results presentation
    \item Recalculation capability
\end{itemize}

\subsection{Adaptive Training}
\begin{itemize}
    \item Personalized task selection
    \item Dynamic difficulty adjustment
    \item Session-based structure
    \item Progress tracking
\end{itemize}

\subsection{Performance Analytics}
\begin{itemize}
    \item Baseline vs current comparison
    \item Radar chart visualization
    \item Domain-specific metrics
    \item Trend analysis
\end{itemize}

\subsection{Gamification}
\begin{itemize}
    \item Badge achievement system
    \item Streak tracking with freeze
    \item Progress milestones
    \item Motivational feedback
\end{itemize}

\subsection{User Experience}
\begin{itemize}
    \item Real-time progress updates
    \item Visual feedback
    \item Intuitive navigation
    \item Responsive design
\end{itemize}

\section{Future Enhancements (Potential)}

\subsection{Advanced Analytics}
\begin{itemize}
    \item Weekly/Monthly progress reports
    \item Performance prediction
    \item Personalized insights
    \item Export data functionality
\end{itemize}

\subsection{Enhanced Gamification}
\begin{itemize}
    \item Leaderboards
    \item Social features
    \item Custom challenges
    \item Reward tiers
\end{itemize}

\subsection{New Task Types}
\begin{itemize}
    \item Additional cognitive domains
    \item Task variations
    \item Custom difficulty curves
    \item Specialized training modules
\end{itemize}

\subsection{Mobile Experience}
\begin{itemize}
    \item Native mobile app
    \item Offline mode
    \item Push notifications
    \item Mobile-optimized tasks
\end{itemize}

\subsection{AI Integration}
\begin{itemize}
    \item Personalized recommendations
    \item Adaptive pacing
    \item Predictive difficulty
    \item Smart scheduling
\end{itemize}

\section{System Status}

\textbf{Current Version:} 1.0  
\textbf{Status:} Production Ready

\subsection{Core Features Implemented}
\begin{itemize}
    \item User authentication
    \item Baseline assessment (6 domains)
    \item Training plan generation
    \item Adaptive difficulty system
    \item Session tracking (4-task sessions)
    \item Streak tracking with freeze
    \item Badge system
    \item Performance comparison
    \item Visual analytics
    \item Cache-busting for fresh data
    \item Real-time progress updates
    \item Difficulty visualization
\end{itemize}

\subsection{Bug Fixes Applied}
\begin{itemize}
    \item Task completion tracking (commit order)
    \item Session counting (4 tasks = 1 session)
    \item Baseline recalculation date update
    \item Dashboard metric accuracy
    \item Streak freeze logic
    \item Cache invalidation
\end{itemize}

\section{Conclusion}

NeuroBloom provides a comprehensive, scientifically-informed cognitive training platform with:

\begin{itemize}
    \item \textbf{Personalized training} based on individual baselines
    \item \textbf{Adaptive difficulty} that grows with user performance
    \item \textbf{Rich analytics} to track improvement
    \item \textbf{Gamification} to maintain engagement
    \item \textbf{Robust tracking} of sessions and streaks
\end{itemize}

The system successfully combines cognitive science principles with modern web development to create an engaging, effective training experience.

\end{document}
